\documentclass[11pt, a4paper]{article}

\usepackage{fontspec}
\setmainfont{Latin Modern Roman}
\setmonofont{DejaVu Sans Mono}[Scale=0.85]
\usepackage[margin=2.5cm]{geometry}
\usepackage{amsmath, amssymb, amsthm}
\usepackage{hyperref}
\usepackage[dvipsnames]{xcolor}
\usepackage{enumitem}
\usepackage{fancyvrb}
\usepackage{caption}
\usepackage{draftwatermark}
\SetWatermarkText{PREPRINT}
\SetWatermarkScale{1.1}
\SetWatermarkLightness{0.92}
\captionsetup[figure]{name=Listing}
\newcommand{\lean}[1]{\ifmmode\text{\ttfamily\detokenize{#1}}\else{\ttfamily\detokenize{#1}}\fi}

\definecolor{indigo}{rgb}{0.29, 0.0, 0.51}
\definecolor{teal}{rgb}{0.0, 0.5, 0.5}
\hypersetup{colorlinks=true, linkcolor=teal, citecolor=indigo, urlcolor=blue,
  pdftitle={Eta-quotients in Lean 4 (DRAFT)}, pdfauthor={Xavier Callens}}

\newtheorem{theorem}{Theorem}[section]
\newtheorem{lemma}[theorem]{Lemma}
\newtheorem{definition}[theorem]{Definition}
\theoremstyle{remark}
\newtheorem{remark}[theorem]{Remark}

\title{\textbf{Eta-Quotients in Lean 4}\\[0.3em]
  \large\textbf{Cusp Orders at Every Cusp, Ligozat's Criterion at Small Level,\\
  and a Named Obstruction at General Level}\\[0.8em]
  \large\textcolor{BrickRed}{PREPRINT --- not peer reviewed}}
\author{
  Xavier Callens\thanks{Corresponding author: \texttt{callensxavier@gmail.com}} \\
  \small SocrateAI Research Initiative
}
\date{\today\\[0.6em]
  \small Archived at Zenodo: \href{https://doi.org/10.5281/zenodo.22648098}{\texttt{10.5281/zenodo.22648098}} (all versions)}

\begin{document}
\maketitle

\begin{center}
\fbox{\begin{minipage}{0.92\textwidth}
\small \textbf{Status.} This is a \emph{preprint}. It has \textbf{not been peer reviewed}. The
Lean development it describes compiles and is axiom-audited --- those claims are machine-checked and
independently reproducible --- but the \emph{exposition} has had no external referee. One headline
result is deliberately bounded (\S\ref{sec:ligozat}): Ligozat's transformation law is proved for
$N \in \{1,2,3,4,5,7,13\}$ and \textbf{not} in general, the general case being blocked by the
obstruction named in \S\ref{sec:obstruction}. Read the scope before citing the result.
\end{minipage}}
\end{center}

\begin{abstract}
We formalize in Lean~4, on top of Mathlib, the arithmetic and analytic layers of Ligozat's
criterion for eta-quotients $f(z) = \prod_{\delta \mid N}\eta(\delta z)^{r_\delta}$, and we report
one theorem, one bounded theorem, and one obstruction. \textbf{(i)} For \emph{every} level $N$,
every integer exponent vector $r$ with $\sum_{\delta\mid N} r_\delta = 2k$, and \emph{every}
$\gamma \in SL_2(\mathbb{Z})$ --- with no congruence condition and no $\Gamma_0(N)$-membership
hypothesis --- we prove the two-sided asymptotic
$\lVert (f\vert_k\gamma)(z)\rVert = \Theta\!\left(\exp(-2\pi\,\mathrm{ord}^{\mathrm{raw}}(N,r,c)\,\mathrm{Im}\,z)\right)$
along $\mathrm{Im}\,z\to\infty$, where $c = \gamma_{10}$ and
$\mathrm{ord}^{\mathrm{raw}}(N,r,c) = \tfrac1{24}\sum_{\delta\mid N}\gcd(c,\delta)^2 r_\delta/\delta$.
At the cusp $\infty$ we additionally identify the order in \emph{Mathlib's own} language,
\lean{meromorphicOrderAt (cuspFunction 1 f) 0}. \textbf{(ii)} We prove Ligozat's transformation law
in full --- $f(\gamma z) = w(\gamma)(cz+d)^k f(z)$ on \emph{all} of $\Gamma_0(N)$, with the
character identified --- for $N \in \{1,2,3,4,5,7,13\}$. \textbf{(iii)} For general $N$ we do not
prove it, and we explain precisely why: the multiplier is
$w(\gamma) = \exp\!\left(\tfrac{\pi i}{12}\mathrm{per}(\gamma)\right)$ where $\mathrm{per}$ is the
period cocycle of the weight-two quasi-modular combination $\sum_\delta r_\delta\,\delta\,E_2(\delta z)$;
for a single $\eta$ this period function is Rademacher's $\Phi$, whose non-coboundary content is the
Dedekind sum $s(d,|c|)$. Mathlib contains no Dedekind sums. We state the minimal Mathlib addition
that would remove the obstruction. The development is 519 declarations in 8117 lines across five
modules, with no \lean{sorry}, and 391 build-failing axiom guards.
\end{abstract}

\noindent\textbf{MSC 2020:} 11F20, 11F11, 68V20.\quad
\textbf{Keywords:} eta-quotients, Ligozat's criterion, Dedekind eta, Dedekind sums,
Rademacher $\Phi$, Lean 4, Mathlib.

\section{Introduction}

An \emph{eta-quotient} of level $N$ is a product
\[
  f(z) \;=\; \prod_{\delta \mid N} \eta(\delta z)^{r_\delta}, \qquad r_\delta \in \mathbb{Z},
\]
over the divisors of $N$, where $\eta$ is Dedekind's function. Ligozat's criterion states that $f$
is a modular form on $\Gamma_0(N)$ of weight $k = \tfrac12\sum_\delta r_\delta$ provided
\begin{equation}\label{eq:ligozat}
  \sum_{\delta\mid N}\delta\, r_\delta \equiv 0 \!\!\pmod{24}, \qquad
  \sum_{\delta\mid N}\tfrac{N}{\delta}\, r_\delta \equiv 0 \!\!\pmod{24},
\end{equation}
and the order of $f$ at every cusp is non-negative~\cite{Ligozat1975, Martin1996, Ono2004}. It is
the standard working tool for producing modular forms by hand, and the reason eta-quotients appear
throughout the theory of partitions and of monstrous moonshine.

This paper reports what happens when one tries to formalize it. The honest answer is: the
arithmetic and the analysis go through in general; the \emph{group theory} does not, and the
obstruction is sharp and classical.

\subsection*{What Mathlib provides, and what it does not}

We are explicit about the substrate because it determines what is new. At commit
\texttt{905b9581} (2026-07-28) Mathlib contains, in \lean{NumberTheory/ModularForms/}:

\begin{itemize}
  \item \lean{DedekindEta.lean} --- 20 declarations: $\eta$ as the $q$-product
        $\eta(z) = q^{1/24}\prod_n(1-q^{n+1})$, non-vanishing (\lean{eta_ne_zero}),
        differentiability, and $\mathrm{logDeriv}\,\eta = \tfrac{\pi i}{12}E_2$;
  \item \lean{Discriminant.lean} --- crucially, the \emph{$S$-transformation of $\eta$}:
        \lean{eta_comp_eq_csqrt_I_inv} gives $\eta(-1/z) = (\sqrt{i})^{-1}\sqrt{z}\,\eta(z)$,
        together with the modularity of $\Delta = \eta^{24}$;
  \item \lean{Cusps.lean} (\lean{IsCusp}, \lean{CuspOrbits}, \lean{widthInfty}),
        \lean{QExpansion.lean} (\lean{cuspFunction}, \lean{qExpansion}), and the general order
        machinery \lean{analyticOrderAt} / \lean{meromorphicOrderAt}.
\end{itemize}

\noindent Repository-wide searches return \textbf{zero} occurrences of \texttt{EtaQuotient},
\texttt{Ligozat}, \texttt{DedekindSum}, \texttt{dedekindSum} and \texttt{etaMultiplier}. The last
two absences are the subject of \S\ref{sec:obstruction}.

\begin{remark}[A methodological warning we paid for]
An earlier pass in this project concluded that Mathlib had no order-of-vanishing machinery, on the
strength of \texttt{grep -rl 'orderAt'} returning zero files. It does have it: the declarations are
\lean{analyticOrderAt} and \lean{meromorphicOrderAt}, and \texttt{grep -rli 'orderat'} returns
nineteen. A zero result from a case-sensitive search on a lower-camelCase name is not evidence of
absence. Correcting this turned an assumed obstruction into Theorem~\ref{thm:infty}.
\end{remark}

\section{The arithmetic layer}

The combinatorial content is separated into a module mentioning no modular form, no $\eta$ and no
upper half-plane: it is a statement about $r$ and $N$ alone. The central definition is Ligozat's
cusp-order expression, for $d \mid N$,
\[
  \mathrm{ord}(N,r,d) \;=\; \frac{N}{24}\sum_{\delta\mid N}
     \frac{\gcd(d,\delta)^2\, r_\delta}{\gcd\!\left(d,\tfrac{N}{d}\right)\, d\, \delta},
\]
transcribed from~\cite[Thm.\ 1.64]{Ono2004}, together with the two congruences
of~\eqref{eq:ligozat} as predicates \lean{LigozatCongr1} and \lean{LigozatCongr2}. The divisor
involution $\delta\mapsto N/\delta$ (Mathlib's \lean{Nat.sum_div_divisors}) gives the duality
exchanging the two congruences, formalized as \lean{ligozatCongr1_dual}.

\section{Cusp orders, and a correction to our own specification}\label{sec:cusp}

\begin{theorem}[Order at every cusp, general $N$]\label{thm:theta}
\lean{etaQuotient_cusp_order}. For every $N$, every $r$ with $\sum_\delta r_\delta = 2k$, and
every $\gamma\in SL_2(\mathbb{Z})$, writing $c = \gamma_{10}$,
\[
  \bigl\lVert (f\vert_k\gamma)(z)\bigr\rVert
  \;=\;\Theta_{\,\mathrm{Im}\,z\to\infty}\;
  \exp\!\bigl(-2\pi\,\mathrm{ord}^{\mathrm{raw}}(N,r,c)\,\mathrm{Im}\,z\bigr),
  \qquad
  \mathrm{ord}^{\mathrm{raw}}(N,r,c) = \frac1{24}\sum_{\delta\mid N}\frac{\gcd(c,\delta)^2 r_\delta}{\delta}.
\]
\end{theorem}

Two features deserve emphasis. There is \emph{no} congruence hypothesis and \emph{no}
$\Gamma_0(N)$-membership hypothesis: the asymptotic holds for every $\gamma$ in the full modular
group. And the $\Theta$ is taken along \lean{atImInfty}, which constrains only $\mathrm{Im}\,z$;
the statement is therefore uniform in $\mathrm{Re}\,z$ and no approach region is assumed.

\begin{remark}[The width normalisation --- our specification was wrong]\label{rem:width}
This node was originally specified with the exponent written as Ligozat's $\mathrm{ord}(N,r,c)$.
\emph{That statement is false}, and it is not what was proved. Ligozat's $\mathrm{ord}(N,r,d)$ is an
order of vanishing in the local uniformiser $q_h = e^{2\pi i z/h}$ at a cusp of denominator $d$,
where $h = N/\gcd(d^2,N)$ is the width of that cusp for $\Gamma_0(N)$. A decay statement in
$\mathrm{Im}\,z$ therefore carries the exponent $\mathrm{ord}/h$, not $\mathrm{ord}$. Formally,
\lean{etaQuotientCuspOrder_eq_width_mul_raw} gives
$\mathrm{ord}(N,r,d) = h\cdot\mathrm{ord}^{\mathrm{raw}}(N,r,d)$, so the two exponentials differ
whenever $h\neq 1$. The companion statement \lean{etaQuotient_cusp_order_ligozat} spells the
exponent as $\mathrm{ord}/h$ so that the correction is visible in the statement itself.

Ligozat's condition (iii) is untouched: $h>0$, so $\mathrm{ord}^{\mathrm{raw}}\ge 0$ at a cusp if
and only if $\mathrm{ord}\ge 0$ there. A hand-computed pin is kept in the development:
$f = (\eta(z)\eta(2z))^{12}$ satisfies $\lVert f\vert_{12}S\rVert \asymp \exp(-3\pi y/2)$.
\end{remark}

\begin{theorem}[The cusp $\infty$, in Mathlib's own language]\label{thm:infty}
\lean{meromorphicOrderAt_cuspFunction_eq_cuspOrder_infty}. For every $N$, every $r$, and every
$m\in\mathbb{Z}$ with $24m = \sum_{\delta\mid N}\delta\,r_\delta$,
\[
  \mathrm{meromorphicOrderAt}\bigl(\mathrm{cuspFunction}_1(f)\bigr)(0) \;=\; m,
\]
and $m = \mathrm{ord}(N,r,N)$, Ligozat's value at $d=N$.
\end{theorem}

This uses Mathlib's own notion of order, its own cusp function, and its own width
($\lean{strictWidthInfty}\,\Gamma_0(N) = 1$); we introduce no competing notion of ``order at a
cusp''. There is no positivity and no holomorphy hypothesis --- the order may be negative.

\section{Ligozat's transformation law at small level}\label{sec:ligozat}

\begin{theorem}[Full transformation law, $N\le 4$]\label{thm:le4}
\lean{ligozat_of_le_four}. Let $0<N\le 4$, let $\sum_{\delta\mid N} r_\delta = 2k$, and assume both
congruences~\eqref{eq:ligozat}. Let $w : \Gamma_0(N)\to\mathbb{C}^\times$ be any homomorphism with
$w(T)=1$, $w(V)=1$ and $w(-I)=(-1)^k$, where $V = W_N T W_N^{-1}$. Then for every
$\gamma\in\Gamma_0(N)$ and every $z\in\mathbb{H}$,
\[
  f(\gamma z) \;=\; w(\gamma)\,(cz+d)^k\, f(z).
\]
\end{theorem}

The architecture is worth stating because it is what makes the bounded result possible. A general
construction produces \emph{a} character with this transformation law for every $N$ and with no
congruence conditions; what it cannot do is say \emph{which} character. Ligozat's first congruence
pins it at $T$; the second pins it at $V = W_N T W_N^{-1}$ --- and $W_N$ here is the Fricke
involution from the companion development~\cite{Fricke2026}, which is where that work is reused; a
separate argument pins it at $-I$. For $N\le 4$ those three elements generate $\Gamma_0(N)$, by a
Euclidean descent, so the character is determined and the congruences are upgraded from
``conditions at two generators'' to the transformation law on the whole group.

\begin{theorem}[Genus-zero primes]\label{thm:primes}
\lean{ligozat_of_prime}. The same conclusion holds for $p\in\{5,7,13\}$.
\end{theorem}

The descent fails at $N=5$, and not by accident: the development records the explicit
counterexample \lean{descent_bound_fails_at_five}. Classically $\overline{\Gamma}_0(5)$ is the free
product $\mathbb{Z}/2 * \mathbb{Z}/2 * \mathbb{Z}$, so $\{-I,T,V\}$ cannot generate it and two
elliptic generators are needed. That necessity is a literature fact and is \emph{not} formalized
here; what is formalized is that this proof needs them --- deleting the two elliptic generators and
re-running the block leaves exactly two unprovable goals. The replacement tool is a Schreier
transversal criterion, \lean{Gamma0_eq_of_schreier}, valid at every prime $p$: a subgroup
containing $-I$, $T$, $V$ equals $\Gamma_0(p)$ as soon as it contains one Schreier generator for
each $j\in[1,p-1]$. We supply those generators at the three primes where $X_0(p)$ has genus zero.

\section{The obstruction at general level}\label{sec:obstruction}

We do not prove Ligozat's transformation law for general $N$, and the reason is not a shortage of
effort or a choice of proof strategy. It is exact.

Write the multiplier as $w(\gamma) = \exp\!\left(\tfrac{\pi i}{12}\,\mathrm{per}(\gamma)\right)$,
where $\mathrm{per}$ is the period cocycle of the weight-two quasi-modular combination
$\sum_{\delta\mid N} r_\delta\,\delta\,E_2(\delta z)$. This period is precisely the integration
constant that both natural soft routes destroy: the $\mathrm{logDeriv}$ route differentiates it
away, and the $24$th-power route (passing to $f^{24}$, a quotient of discriminants whose modularity
Mathlib already has) quotients it away. For a single $\eta$ this period function \emph{is}
Rademacher's $\Phi$,
\[
  \Phi(\gamma) \;=\; \frac{a+d}{c} \;-\; 12\,\mathrm{sign}(c)\, s(d,|c|)
     \;-\; 3\,\mathrm{sign}\bigl(c(a+d)\bigr),
\]
whose non-coboundary content is exactly the Dedekind sum $s(d,|c|)$.

For any level with $\mathrm{genus}(X_0(N))\ge 1$ --- $N=11$ and most $N\ge 17$ --- the group
$\Gamma_0(N)$ has hyperbolic generators, and no soft argument evaluates $w$ there: on the free part
of $\Gamma_0(N)^{\mathrm{ab}}$, $w$ is constrained only by $w^{24}=1$, and
$\mathrm{Hom}(\Gamma_0(N),\mu_{24})$ is large. Evaluating $w$ in closed form on a single hyperbolic
element of $\Gamma_0(11)$ is mathematically equivalent to evaluating a Dedekind sum. Every disguise
we attempted reduces back to this: Atkin--Lehner conjugation reintroduces general eta transforms,
theta/Poisson routes re-derive the same $24$th-root phases, and Wohlfahrt level-$24$ arguments
presuppose the multiplier formula.

\subsection*{What would remove it}

In ascending strength, the minimal Mathlib additions are:
\begin{enumerate}[label=(\arabic*)]
  \item the Dedekind sum $s(d,c)$ and reciprocity;
  \item Rademacher's $\Phi : SL_2(\mathbb{Z})\to\mathbb{Z}$ with the cocycle law
        $\Phi(AB) = \Phi(A)+\Phi(B)-3\,\mathrm{sign}(c_A c_B c_{AB})$;
  \item the Petersson--Rademacher multiplier
        $\eta(\gamma z) = \exp\!\left(\pi i\left(\tfrac{a+d}{12c} - s(d,c) - \tfrac14\right)\right)\sqrt{cz+d}\;\eta(z)$
        for $c>0$.
\end{enumerate}
With (3), general-$N$ Ligozat becomes moderate bookkeeping on top of the argument already
formalized here. We note that (2) is best constructed as the period of $E_2$ using the
\lean{E2_slash_action} machinery \emph{already present} in Mathlib --- the analytic half exists ---
which makes it the cleanest available upstream contribution.

\begin{remark}[A per-instance escape hatch]
For a \emph{fixed} $N$ and $r$, the multiplier on a finite generating set is a ratio of convergent
$q$-products valued in the discrete set $\mu_{24}$. Interval arithmetic with error below
$|1-\zeta_{24}|/2$ therefore \emph{proves} each value. This is an $A2$ (program-checked) route for
any concrete eta-quotient at any level, requiring no Mathlib addition, though heavy to promote to
kernel-checked $A1$.
\end{remark}

\section{Verification status}

\lean{lake build SocrateAI} completes in 3462 jobs with no errors. The development is
519 declarations in 8117 lines across five modules
(\lean{EtaQuotient}, \lean{EtaQuotientCuspOrder}, \lean{EtaQuotientCuspTheta},
\lean{EtaQuotientModularity}, \lean{EtaQuotientPrimeLevel}). There are no \lean{sorry} placeholders
and no project-local axioms. Following~\cite{FLT2026}, the axiom footprint is enforced rather than
reported: \lean{FinalCheck.lean} carries 391 \texttt{\#guard\_msgs in \#print axioms} guards over
389 distinct theorems, so a widened footprint is a compile error; 370 report exactly
$[\texttt{propext},\ \texttt{Classical.choice},\ \texttt{Quot.sound}]$. A negative control asserting
a deliberately wrong footprint is kept outside the build and verified to fail, which is what
establishes the guards are load-bearing rather than vacuous.

Of the 453 theorems in the new modules, 14 --- three percent --- have a one-line
\lean{rfl}/\lean{decide} proof, and 45\% quantify over a structure. We report this ratio because an
earlier module in this project consisted of 174 numeral identities of which 158 were one-line, and
carried no evidential weight; the check is now standard practice here for any large batch.

The statement-dependency graph (\lean{dag/theorems.jsonl}, 57 nodes, 51 proved, acyclic) is
machine-validated: every node marked proved must name a declaration that resolves and may depend
only on proved nodes. \textbf{The node \lean{ETA-01}, Ligozat's criterion in general, remains
\lean{open} with \lean{lean_name: null}}, and the obstruction node is \lean{blocked} with
\lean{lean_name: null} --- for that node the absence of a declaration \emph{is} the content.

\section{Related work and prior art}

Mathlib has no eta-quotients. The Fermat's Last Theorem artifact~\cite{FLT2026} contains
substantial eta-related material, but a direct reading shows it is about modular \emph{units} on
modular function fields --- \lean{modularUnitSeries}, \lean{sharpUnit*}, over Laurent and Hahn
series. Its three occurrences of ``Ligozat'' are \emph{namespace labels}
(\lean{LigozatUnitEngine}, \lean{LigozatUnitAL}) on that machinery, not a formalization of the
modularity criterion. A GitHub-wide \texttt{language:lean} search returns zero for
\texttt{Ligozat}. We did \emph{not} search the Lean Zulip or open mathlib4 pull requests, and we do
not claim coverage there.

\section{Limitations}

\begin{remark}[What is not done]
\leavevmode
\begin{enumerate}[label=(\alph*)]
  \item Ligozat's criterion for general $N$ is not proved (\S\ref{sec:obstruction}).
  \item The order of vanishing at a cusp \emph{other than} $\infty$ is not stated in Mathlib's
        \lean{meromorphicOrderAt} language. Mathlib's width is \lean{widthInfty} only; there is no
        \lean{widthAtCusp} and no \lean{orderAtCusp}. These are constructible from \lean{conjGL},
        \lean{widthInfty} and the slash action, so this is missing API rather than an obstruction.
        Theorem~\ref{thm:theta} gives the analytic content at every cusp regardless.
  \item The indexing data is not formalized: that the cusps of $\Gamma_0(N)$ are parametrised by
        $d\mid N$ with multiplicity $\varphi(\gcd(d,N/d))$, and that $[SL_2(\mathbb{Z}):\Gamma_0(N)]
        = \psi(N)$. No index formula for $\Gamma_0$ exists anywhere in Mathlib. The definitions
        \lean{numCusps} and \lean{dedekindPsi} in this development carry those classical readings
        but do not prove them.
  \item Theorems~\ref{thm:le4} and~\ref{thm:primes} give the transformation law; assembling it with
        the cusp-order positivity condition into membership in a bundled \lean{ModularForm} is not
        yet done.
  \item That $\{-I,T,V\}$ \emph{cannot} generate $\Gamma_0(5)$ is cited, not proved.
\end{enumerate}
\end{remark}

\begin{thebibliography}{99}
\bibitem{Ligozat1975} G.~Ligozat, \textit{Courbes modulaires de genre 1}, Bull.\ Soc.\ Math.\ France, M\'em.\ \textbf{43} (1975), 1--80.
\bibitem{Martin1996} Y.~Martin, \textit{Multiplicative eta-quotients}, Trans.\ Amer.\ Math.\ Soc.\ \textbf{348} (1996), 4825--4856.
\bibitem{Ono2004} K.~Ono, \textit{The Web of Modularity}, CBMS Regional Conference Series in Mathematics \textbf{102}, AMS (2004). \texttt{doi:10.1090/cbms/102}.
\bibitem{Rademacher1973} H.~Rademacher, \textit{Topics in Analytic Number Theory}, Grundlehren \textbf{169}, Springer (1973).
\bibitem{Apostol1990} T.~M.~Apostol, \textit{Modular Functions and Dirichlet Series in Number Theory}, 2nd ed., GTM \textbf{41}, Springer (1990).
\bibitem{Mathlib2020} The mathlib Community, \textit{The Lean mathematical library}, CPP 2020.
\bibitem{FLT2026} Anthropic, \textit{A formal proof of Fermat's Last Theorem in Lean 4}, software artifact, \url{https://github.com/anthropics/fermats-last-theorem}, 2026.
\bibitem{Fricke2026} X.~Callens, \textit{A Mathlib-Idiomatic Formalization of the Fricke Involution and its Operator on Modular Forms}, Zenodo, 2026. \texttt{doi:10.5281/zenodo.22542571}.
\bibitem{ThisWork} X.~Callens, \textit{Eta-Quotients in Lean 4} (this preprint), Zenodo, 2026. \texttt{doi:10.5281/zenodo.22648098}.
\bibitem{deMoura2021} L.~de~Moura et al., \textit{The Lean 4 Theorem Prover and Programming Language}, CADE-28, LNCS \textbf{12699}, Springer (2021).
\end{thebibliography}

\end{document}
